% high_moment_inertia.tex -- round 398, reviewer-preregistered
% high-moment kill-test of P1 = ind_-(R_{N-3}-I/2) <= 1,
% before the source-Weyl campaign.  Outcome: KILL_FAIL.
% Builds with pdflatex.
% Companion machine check: verify_high_moment.py.
% Discovery census: experiments/tfpt-discovery/high_moment_inertia_probe.py.
% NO RH CLAIM.  Research documentation, not a theorem of RH.
\documentclass[11pt]{article}

\usepackage[a4paper,margin=2.7cm]{geometry}
\usepackage{amsmath,amssymb,amsthm}
\usepackage{booktabs}
\usepackage{microtype}
\usepackage{xcolor}
\usepackage[colorlinks=true,linkcolor=blue!50!black,urlcolor=blue!50!black]{hyperref}

\newtheorem{lemma}{Lemma}
\newtheorem{prop}{Proposition}
\newtheorem{definition}{Definition}
\theoremstyle{remark}
\newtheorem{remark}{Remark}

\newcommand{\R}{\mathbb{R}}
\newcommand{\N}{\mathbb{N}}
\newcommand{\indm}{\operatorname{ind}_{-}}
\newcommand{\half}{\tfrac12}
\newcommand{\tr}{\operatorname{tr}}

\title{\bfseries High-moment inertia, killed\\[2mm]
\large A fixed even cycle-moment does not certify
$\indm(R_{N-3}-\tfrac12 I)\le 1$}
\author{Stefan Hamann}
\date{August 29, 2026}

\begin{document}
\maketitle

\begin{abstract}\noindent
The dual resolvent contraction $R_{N-3}$ of a canonical window
is a positive contraction.
For every integer $d\ge 1$ the pointwise bound
$1_{\{r<1/2\}}\le [2(1-r)]^{2d}$ yields
$\indm(R-\tfrac12 I)\le 2^{2d}\tr((I-R)^{2d})$.
If the right-hand side were $<2$ for one \emph{fixed}
$d\in\{2,3,4,6,8\}$ on MAIN, Twin and living $\chi$, with
reserve growing in $N$, then P1 would be a finite source
cycle-sum, outside the two-moment ceiling class.
It is not.
On the named census the moment is $\ge 16$ already at $d=8$
on FRAME-A, grows with $N$ (core-$42$: $M_2$ from $34.9$ at
$N=142$ to $224.5$ at $N=878$), and no $d$ --- even $d>8$ ---
brings it under $2$ on the $n_{-}=1$ branch.
The cause is a macroscopic cluster of eigenvalues of $R_{N-3}$
at $\tfrac12$ from above (FRAME-A: $36$ in $[\tfrac12,\tfrac35]$),
each contributing $\sim 1$ to the sum.
Scramble sees $\indm=21$; two-period dies; dead $\chi$ is
$\ge 2$ as a control, but so is living $\chi$ --- the test is
not a world separator.
No RH claim.
\end{abstract}

\medskip\noindent
\textbf{Status.}\enspace
Lemma~\ref{lem:bound} and Lemma~\ref{lem:cycle} are proved
(finite identities).
Proposition~\ref{prop:fail} names the kill.
No RH claim.  No $L^*$ claim.  No $R^{\dagger}$ claim.

\section{The object}\label{sec:obj}

Let $R_{N-3}$ be the dual Christoffel--Darboux Gram of the first
$N-3$ dual orthogonal polynomials, restricted to the hole nodes
$Y$ --- the rest block of the r367 two-rank cut~\cite{r367,r369}.
On a window $(\mu,\nu)$ with half-filling $S=2N-1$,
\[
R_{N-3}
\;=\;
B_{\bullet,\,0:N-3}\,B_{\bullet,\,0:N-3}^{T},
\]
$B$ the dual OP matrix dressed by $\sqrt{u^{\vee}}$ and
evaluated on $Y$.
Haynsworth reduces window-local $R^{\dagger}\succ \half I$ to
(P1) $\indm(R_{N-3}-\half I)\le 1$ and an $O(1)$ edge
sign~\cite{rdag}.
The present test is a \emph{sufficient} majorant for (P1),
preregistered before any source-Weyl energy estimate.

Hard stops, binding: $d$ is chosen from the five-element set
$\{2,3,4,6,8\}$, never a function of $N$; the certificate is a
trace of an even power, not an eigenvector; no target margin
enters.

\section{The bound}\label{sec:bound}

\begin{lemma}[even moment majorant]\label{lem:bound}
Let $0\preceq R\preceq I$ on a finite-dimensional real Hilbert
space, with eigenvalues $r$.
For every integer $d\ge 1$ and every $r\in[0,1]$,
\[
1_{\{r<1/2\}}
\;\le\;
\bigl[2(1-r)\bigr]^{2d}.
\]
Summing,
$\indm(R-\half I)\le 2^{2d}\tr\bigl((I-R)^{2d}\bigr)$.
\end{lemma}

If $r<\half$ then $2(1-r)>1$, so the even power is $\ge 1$.
If $r\ge\half$ the indicator vanishes and the right-hand side
is nonnegative.
Over $\mathbb{Q}$: $r=\tfrac25$, $d=2$ gives
$[2(1-\tfrac25)]^{4}=(\tfrac65)^{4}=\tfrac{1296}{625}\ge 1$.
Dropping the prefactor $2^{2d}$ makes the bound false:
$(1-\tfrac25)^{4}=\tfrac{81}{625}<1$.
An odd power is not a majorant without the contraction
hypothesis: $R=\operatorname{diag}(\tfrac25,2)$ has
$\indm=1$ but
$\sum_i[2(1-r_i)]^{5}<0$.

\section{Cycle sums}\label{sec:cycle}

\begin{lemma}[closed walks]\label{lem:cycle}
For a real symmetric $A$,
\[
\tr(A^{k})
\;=\;
\sum_{i_1,\ldots,i_k}
A_{i_1 i_2}\cdots A_{i_k i_1},
\]
the sum of weighted closed walks of length $k$ on the
complete graph with edge weights $A_{ij}$.
On the dual CD graph, $A=I-R_{N-3}$ and $k=2d$ even.
\end{lemma}

Exact on a $2\times 2$ toy: $A=\bigl(\begin{smallmatrix}1/2&1/4\\1/4&1/3\end{smallmatrix}\bigr)$,
$\tr(A^{2})=\tfrac{35}{72}$.
On a $3\times 3$ numerical toy the $k=4$ walk-sum matches
$\tr(A^{4})$ to $10^{-17}$.
Because the kill-test fails, this identity is \emph{not} the
R399 target: a source expansion of $\tr((I-R)^{2d})$ would
bound a quantity that is not $<2$.

\section{The census, KILL\_FAIL}\label{sec:census}

\begin{prop}[KILL\_FAIL]\label{prop:fail}
On every named living window, for every
$d\in\{2,3,4,6,8\}$,
$M_d:=2^{2d}\tr((I-R_{N-3})^{2d})\ge 2$.
The reserve $2-M_d$ shrinks as $N$ grows.
No registered $d$ therefore reduces P1 to a cycle-moment
theorem.
\end{prop}

FRAME-A ($w=9$, $N=184$, $|Y|=104$): $0\preceq R_{N-3}\preceq I$
(gate), $\indm=1$, $r_{\min}=0.46153$,
\[
M_{2},M_{3},M_{4},M_{6},M_{8}
\;=\;
33.48,\;26.79,\;22.75,\;18.23,\;16.03.
\]
Eigenvalues in $[\half-\delta,\half+\delta]$:
$5$, $24$, $37$ for $\delta=0.01,0.05,0.1$;
$36$ lie in $[\half,\tfrac35]$.
The r367 rest $A_0+\half I$ agrees with $R_{N-3}$ at residual $0$.

Core-$42$: $42/42$ contractions, $\indm\in\{0,1\}$,
$M_2\in[33.22,224.5]$.
At smallest $N=142$ ($kz=18$), $M_2=34.86$;
at largest $N=878$ ($kz=52$), $M_2=224.5$ --- reserve shrinks.
The MAIN-$85$ sample continues the growth (kz$119$: $M_2=287.9$;
EXT $kz=69$, $N=5690$: $M_2=1079$).
Twin $kz\in\{9,18,52\}$: $|M_2^{\mathrm{twin}}-M_2^{\mathrm{MAIN}}|\le 2\cdot 10^{-5}$,
same inertia.

Living $\chi_3$ ($37$ windows) $M_2\in[28.03,183.4]$;
living $\chi_4$ ($41$ windows) $M_2\in[26.90,178.7]$;
all $\ge 2$.
Dead sprouts $\chi_3\in\{15,19,23,33,39\}$ and $\chi_4=20$ are
likewise $\ge 2$ (the control holds) --- but living worlds fail
too, so the test is not sharp as a world discriminator.

\section{Why it dies, and why no larger $d$ saves it}\label{sec:why}

Write $M_d=\sum_i[2(1-r_i)]^{2d}$.
Modes with $r>\half$ have $2(1-r)<1$ and decay in $d$;
a mode with $r<\half$ has $2(1-r)>1$ and \emph{grows} in $d$.
On FRAME-A the minimum over $d\in\mathbb{N}$ is $M_{11}=15.06$,
still $\gg 2$, then $M_{30}=90$.
On $kz=18$ ($r_{\min}=0.390$) the minimum is $M_5=28.66$, then
$M_{30}=1.5\cdot 10^{5}$.
On the vacuous branch $\indm=0$ ($kz=12$, $r_{\min}=0.500018$)
the bound is monotone decreasing, but the cluster sits at
distance $O(10^{-5})$ from $\half$: even the bulk-only sum
drops below $2$ only at $d\sim 1024$, which grows with
$1/\operatorname{gap}$ and is forbidden.

The two-moment ceiling class cannot see this cluster; a
growing even moment cannot either, once $d$ is not allowed
to track $N$.
P1 remains an integer mode count.
The next object is source Weyl energy of the centred mass
difference after the three edge modes, as preregistered, not
a cycle-sum of $I-R_{N-3}$.

\section{Named kills}\label{sec:kills}

Scramble (matched $\chi$-comb, seed $1$): $\indm=21$,
$r_{\min}=0.0036$, $M_2=148$, $M_8=2.65\cdot 10^{5}$.
The trace sees the overload; $M_d$ explodes in $d$.
Two-period ($S=21$, $c=\tfrac23$): $\indm=4$, $r_{\min}=0$,
$M_2=52$.
Mutants: $2^{2d}$ omitted fails over $\mathbb{Q}$; odd
exponent fails on a non-contraction.
Dead $\chi$ is $\ge 2$, as required of a sharp test --- and
uninformative here, because living worlds are $\ge 2$ as well.

\section*{Machine check}

Standalone: Fractions bound, nopref, odd mutant, $2\times 2$
cycle-sum, frozen $d$-set, $3\times 3$ walk-sum.
Construction: FRAME-A contraction and r367 residual,
moments, scramble $21$, two-period, dead $\chi_3$-$15$,
living $\chi_3$-$9$ with $\indm=0$.
Discovery census: \texttt{high\_moment\_inertia\_probe.py}
(\texttt{--smoke} / full), core-$42$, MAIN sample, Twin,
living and dead $\chi$.
Exit line: \texttt{HIGH MOMENT VERIFIED}.

\section*{Claim boundary}

Finite identities for the even-moment majorant and for closed
walks; a named census that the sufficient test of P1 fails on
every registered $d$, with reserve shrinking in $N$; named
scramble / two-period / mutant / dead-$\chi$ controls.
No statement about zeros of $\zeta(s)$, in either direction.
No RH claim.

\begin{thebibliography}{9}
\bibitem{r367}
S.~Hamann,
\textit{Final two-rank inertia},
sealed probe, August 2026,
\texttt{experiments/tfpt-discovery/final\_two\_rank\_inertia\_probe.py}.
\bibitem{r369}
S.~Hamann,
\textit{Mixed Haynsworth},
sealed probe, August 2026,
\texttt{experiments/tfpt-discovery/mixed\_haynsworth\_probe.py}.
\bibitem{rdag}
S.~Hamann,
\textit{The saturated dual resolvent},
standalone note, August 2026,
\texttt{rh/problem/rdagger\_saturation.tex}.
\end{thebibliography}

\end{document}
